\documentclass[10pt,a4paper]{article}
\usepackage{amsmath,amssymb,amsthm,graphicx}
\usepackage{makecell}
% Card-local minimal preamble (journal unit sigma-2025-023-macdonald-identities-weyl-kac-denominator-fo, P2-ahl16 precedent).
%% Verbatim from the pinned source: the author theorem environments and macro block,
%% so every macro and environment of the counted spans is available.
% Omitted and why: the documentclass line (the journal class sigma is not on the host;
%% the assembler supplies the class) and page-geometry/font blocks (they fight the
%% forced xelatex engine).
\usepackage{amsmath,amssymb,amsthm,mathtools}

\renewcommand{\textfraction}{0.01}
\renewcommand{\topfraction}{0.99}
\renewcommand\theadalign{bc}
\renewcommand\theadfont{\bfseries}
\renewcommand\theadgape{\Gape[4pt]}
\renewcommand\cellgape{\Gape[4pt]}
\newtheorem{Theorem}{Theorem}[section]
\newtheorem{Corollary}[Theorem]{Corollary}
\newtheorem{Lemma}[Theorem]{Lemma}
\newtheorem{Proposition}[Theorem]{Proposition}
 { \theoremstyle{definition}
\newtheorem{Definition}[Theorem]{Definition}
\newtheorem{Note}[Theorem]{Note}
\newtheorem{Example}[Theorem]{Example}
\newtheorem{Remark}[Theorem]{Remark} }
\numberwithin{equation}{section}
\newcommand{\ti}{\tilde}
\newcommand{\al}{\alpha}
\newcommand{\om}{\omega}
\newcommand{\ga}{\gamma}
\newcommand{\la}{\lambda}
\newcommand{\eps}{\varepsilon}
\newcommand{\si}{\sigma}
\newcommand{\La}{\Lambda}
\newcommand{\Ga}{\Gamma}
\newcommand{\ov}[1]{\overline{#1}}
\newcommand{\os}[2]{\overset{#1}{#2}}
\newcommand{\scb}[2]{\scalebox{#1}{#2}}
\newcommand{\crossp}{\overset{+}{\leadsto}}
\newcommand{\crossn}{\overset{-}{\leadsto}}
\newcommand{\fh}{\mathfrak{h}}
\newcommand{\fg}{\mathfrak{g}}
%
\newcommand{\cC}{\mathcal{C}}
\newcommand{\cF}{\mathcal{F}}
\newcommand{\cU}{\mathcal{U}}
\newcommand{\cM}{\mathcal{M}}
\newcommand{\cA}{\mathcal{A}}
\newcommand{\cB}{\mathcal{B}}
\def\cch{\mathcal H}
\def\ccsc{\mathcal {SC}}
\def\ccdd{\mathcal {DD}}
\def\ccd{\mathcal {D}}
\def\ccp{\mathcal P}
\def\ccc{\mathcal C}
\DeclareMathOperator{\BGP}{BG}
\DeclareMathOperator{\WLP}{WL}
\DeclareMathOperator{\core}{core}
\DeclareMathOperator{\quot}{quot}
\newcommand{\scal}[2]{\langle #1,#2\rangle}
\newcommand{\scald}[2]{(#1,#2)}
\newcommand{\qu}[1]{\qquad\text{#1}\qquad}
\newcommand{\te}[1]{\text{ #1 }}
\newcommand{\Z}{\mathbb{Z}}
\newcommand{\N}{\mathbb{N}}
\newcommand{\C}{\mathbb{C}}
\newcommand{\R}{\mathbb{R}}
\newcommand{\K}{\mathbb{K}}
\newcommand{\Q}{\mathbb{Q}}
\newcommand{\bB}{\mathbb{B}}
\newcommand{\bA}{\mathbb{A}}
\newcommand{\fB}{\mathfrak{B}}
\newcommand{\wB}{\widehat{B}}
\newcommand{\mk}{\mathrm{k}}
\newcommand{\mK}{\mathrm{K}}
\newcommand{\sLa}{\mathsf{\La}}
\newcommand{\cAJ}{\mathcal{A}_J}
\DeclareMathOperator{\Sp}{Sp}
\DeclareMathOperator{\GL}{GL}
\DeclareMathOperator{\OO}{SO}
\renewcommand{\OE}{\mathrm{O}}
\DeclareMathOperator{\I}{I}
\DeclareMathOperator{\oo}{so}
\renewcommand{\oe}{\mathrm{o}^{\text{even}}}
\newcommand{\func}[1]{\operatorname{#1}}
\newcommand{\lra}{\longrightarrow}
\newcommand{\cev}[1]{\reflectbox{\ensuremath{\vec{\reflectbox{\ensuremath{#1}}}}}}
\newcommand{\cI}{\mathcal{I}}
\newcommand\blfootnote[1]{%
	\begingroup
	\renewcommand\thefootnote{}\footnote{#1}%
	\addtocounter{footnote}{-1}%
	\endgroup
}
\newcommand{\syc}{\mathsf{s_C}}
\newcommand{\lb}{\llbracket}
\newcommand{\rb}{\rrbracket}
\title{Affine Grassmannian elements of type C\_n\textasciicircum{}(1) are parametrized by self-conjugate 2n-cores}
\author{Lecouvey, C\textbackslash{}'\{e\}dric, Wahiche, David}
\date{doi:10.3842/SIGMA.2025.023; CC BY 4.0}
\begin{document}
\maketitle
\setcounter{section}{5}
\setcounter{Theorem}{2}
\noindent\textit{Editorial note.} All delivered mathematical text is the author's own, copied from the pinned journal source file sources/sigma-2025-023-macdonald-identities-weyl-kac-denominator-fo/pinned-original-source.tex (Macdonald Identities, Weyl-Kac Denominator Formulas and Affine Grassmannian Elements); no statement or proof step was rewritten, reordered or removed.\par\medskip
\noindent\textit{Editorial note.} The counted proof is the author's complete argument: The author asserts the stability step ('One first checks') without executing it; the descent crux carries its reason clause ('because c is self-conjugate') and the strictly-decreasing descent terminates, which is implicit. The census printed field 7.1 is the published-version coordinate; the card face renders Lemma 5.3 by source-text counter stepping (census v3.5 rule). No numbered external references are invoked inside the span.\par\medskip
\noindent\textit{Editorial note.} The label Lemma\_selfconjugate names the counted statement, and the proof environment immediately following it in the source is delivered verbatim as its proof.\par\medskip
\noindent\textit{Editorial note.} The counted statement is carried by the author's own Lemma environment in the article source.\par\medskip
\noindent\textit{Editorial note.} Journal-unit assembly: card-local preamble (author theorem environments and macros verbatim), article class, no bibliography in the counted span (the span carries no \textbackslash{}cite commands).\par\medskip
\noindent\textit{Editorial note.} Revision r002: the r001 card-local preamble was malformed - the lane's preamble extractor dropped five structural lines of the author's \textbackslash{}blfootnote definition (\textbackslash{}begingroup, \textbackslash{}addtocounter, \textbackslash{}endgroup, the closing brace, a bare \% comment), leaving an unbalanced brace and an xelatex failure; the extractor is fixed and the preamble is regenerated verbatim from the pinned source. Additionally \textbackslash{}usepackage\{makecell\} is declared via assembly.extra\_packages because the source preamble redefines the journal's tablehead macros (\textbackslash{}theadalign, \textbackslash{}theadfont, \textbackslash{}theadgape, \textbackslash{}cellgape via \textbackslash{}Gape) that the article class does not define; no card-body text changed.\par\medskip
\noindent\textit{Editorial scope.} the counted statement and its own proof environment only\par\medskip
\noindent\textit{Allowed standard prerequisites (editorial).}\begin{itemize}
\item The action of the affine Grassmannian elements W\_\{C\_n\textasciicircum{}(1)\} on core partitions via the raised operators s\_i\textasciicircum{}\{C\_n\textasciicircum{}(1)\} (section 5 setup).
\item The set C\_\{2n\}\textasciicircum{}s of self-conjugate 2n-cores and the notion of removable nodes of a partition.
\end{itemize}\medskip
\begin{Lemma}
	\label{Lemma_selfconjugate}We have $W_{^{C_{n}^{(1)}}}\cdot \varnothing =
	\mathcal{C}_{2n}^{s}$, that is, the affine Grassmannian elements of type $
	C_{n}^{(1)}$ are parametrized by the self-conjugate $2n$-cores.
\end{Lemma}

\begin{proof}
	One first checks that \smash{$\mathcal{C}_{2n}^{s}$} is stable under the action of \smash{$W_{^{C_{n}^{(1)}}}$}.
    Next, it suffices to ob\-serve~that for any nonempty $%
	c\in \mathcal{C}_{2n}^{s}$, there exists at least an integer $i=0,\dots ,n$
	such that all~$i$~and $2n-i$ nodes of $c$ are removable because $c$
	is self-conjugate. This implies that \smash{$ s_{i}^{C_{n}^{(1)}}\cdot c=c^{\prime }$}
    where $c^{\prime }\neq c$ belongs to \smash{$\mathcal{C}_{2n}^{s}$} and has a
	number of nodes strictly less than $c$.
\end{proof}

\end{document}
